% GENERATED FILE. Edit canonical lessons, then run npm run generate:books.
% canonical-source-hash: fnv1a64:824ac79b618349a5
% canonical-lessons: LA-C317-incendium, LA-C317-aestus, LA-C317-favilla, LA-C317-torris, LA-C317-ardor

\chapter{Incendium}
\label{ch:la-incendium}
\hlchaptermodality{317}

\begin{chapteropening}
\textbf{By the end of this chapter:} \emph{I can describe a fire: a blaze, the heat, glowing ash, a firebrand and burning heat.}

Complete the last lesson of chapter 317: I can describe a fire: a blaze, the heat, glowing ash, a firebrand and burning heat.
\end{chapteropening}

\section[incendium, incendiī]{incendium, incendiī — a fire, a blaze}

\label{lesson:LA-C317-incendium}



\begin{quote}
Before the new one: say the Latin for to run to help, then the Latin for help.
\end{quote}

\subsection*{You'll want to know: incendium, incendiī}
\textbf{incendium, incendiī} — \textquotedblleft{}a fire, a blaze\textquotedblright{}.

\begin{quote}
Noctū magnum incendium in urbe ortum est.
\end{quote}

\begin{quote}
\emph{At night a great fire broke out in the city.}
\end{quote}

Neuter, from \emph{incendō}, 'to set on fire' (taught). Fire was the great fear of a crowded city of wooden floors and oil lamps.

\begin{cousinweb}
\textbf{incendium} is built on \emph{incendō}, \textquotedblleft{}to set on fire\textquotedblright{} (EtymDB), from \emph{in} and \emph{candeō}, \textquotedblleft{}to glow\textquotedblright{}. English \emph{incendiary} comes from its family.
\end{cousinweb}

\begin{grammarlens}[title={What you've built}]
One more word for telling, explaining and describing at B1.
\end{grammarlens}

\subsection*{Your turn}
\begin{itemize}
\raggedright
  \item \emph{Say it:} \emph{incendium, incendiī}
  \item \emph{Say it:} \emph{incendium, incendiī}, once more
  \item \emph{Say it:} a sentence of your own with \emph{incendium, incendiī}
  \item \emph{Recall:} say the Latin for to run to help, then the Latin for help, then say \emph{incendium, incendiī} again
\end{itemize}

\begin{tcolorbox}[breakable,colback=teal!4,colframe=teal!35!black,title={Before you move on}]
What does incendium, incendiī mean? (\textquotedblleft{}A fire, a blaze\textquotedblright{}.) Say it once more.
\end{tcolorbox}
\section[aestus, aestūs]{aestus, aestūs — heat, glow; the tide}

\label{lesson:LA-C317-aestus}



\begin{quote}
Before the new one: say the Latin for help, then the Latin for a fire.
\end{quote}

\subsection*{You'll want to know: aestus, aestūs}
\textbf{aestus, aestūs} — \textquotedblleft{}heat, glow; the tide\textquotedblright{}.

\begin{quote}
Propter aestum incendiī nēmō appropinquāre poterat.
\end{quote}

\begin{quote}
\emph{Because of the heat of the fire nobody could come near.}
\end{quote}

Masculine, 4th declension. The heat of fire or summer, the surge of the sea, and the tide. Related to \emph{aestās} (taught).

\begin{cousinweb}
\textbf{aestus} is inherited from a Proto-Indo-European root meaning \textquotedblleft{}to burn\textquotedblright{} (EtymDB), and is a cousin of \emph{aestās}, \textquotedblleft{}summer\textquotedblright{}.
\end{cousinweb}

\begin{grammarlens}[title={What you've built}]
One more word for telling, explaining and describing at B1.
\end{grammarlens}

\subsection*{Your turn}
\begin{itemize}
\raggedright
  \item \emph{Say it:} \emph{aestus, aestūs}
  \item \emph{Say it:} \emph{aestus, aestūs}, once more
  \item \emph{Say it:} a sentence of your own with \emph{aestus, aestūs}
  \item \emph{Recall:} say the Latin for help, then the Latin for a fire, then say \emph{aestus, aestūs} again
\end{itemize}

\begin{tcolorbox}[breakable,colback=teal!4,colframe=teal!35!black,title={Before you move on}]
What does aestus, aestūs mean? (\textquotedblleft{}Heat, glow; the tide\textquotedblright{}.) Say it once more.
\end{tcolorbox}
\section[favilla, favillae]{favilla, favillae — glowing ash, embers}

\label{lesson:LA-C317-favilla}



\begin{quote}
Before the new one: say the Latin for a fire, then the Latin for heat.
\end{quote}

\subsection*{You'll want to know: favilla, favillae}
\textbf{favilla, favillae} — \textquotedblleft{}glowing ash, embers\textquotedblright{}.

\begin{quote}
Ventus favillam in tēcta vīcīna ferēbat.
\end{quote}

\begin{quote}
\emph{The wind was carrying the glowing ash onto the roofs nearby.}
\end{quote}

Feminine, mostly poetic: the hot ash that drifts and starts new fires.

\begin{cousinweb}
\textbf{favilla} is linked by Lewis \& Short with a root meaning \textquotedblleft{}to shine\textquotedblright{}, seen in Greek \emph{phōs}, \textquotedblleft{}light\textquotedblright{}; the link is not certain.
\end{cousinweb}

\begin{grammarlens}[title={What you've built}]
One more word for telling, explaining and describing at B1.
\end{grammarlens}

\subsection*{Your turn}
\begin{itemize}
\raggedright
  \item \emph{Say it:} \emph{favilla, favillae}
  \item \emph{Say it:} \emph{favilla, favillae}, once more
  \item \emph{Say it:} a sentence of your own with \emph{favilla, favillae}
  \item \emph{Recall:} say the Latin for a fire, then the Latin for heat, then say \emph{favilla, favillae} again
\end{itemize}

\begin{tcolorbox}[breakable,colback=teal!4,colframe=teal!35!black,title={Before you move on}]
What does favilla, favillae mean? (\textquotedblleft{}Glowing ash, embers\textquotedblright{}.) Say it once more.
\end{tcolorbox}
\section[torris, torris]{torris, torris — a firebrand, a burning log}

\label{lesson:LA-C317-torris}



\begin{quote}
Before the new one: say the Latin for heat, then the Latin for glowing ash.
\end{quote}

\subsection*{You'll want to know: torris, torris}
\textbf{torris, torris} — \textquotedblleft{}a firebrand, a burning log\textquotedblright{}.

\begin{quote}
Servus torrem ē focō rapuit.
\end{quote}

\begin{quote}
\emph{A slave snatched a burning log from the hearth.}
\end{quote}

Masculine, 3rd declension: a piece of wood still burning.

\begin{cousinweb}
\textbf{torris} belongs with \emph{torreō}, \textquotedblleft{}to scorch, to roast\textquotedblright{} (Lewis \& Short): the scorching log.
\end{cousinweb}

\begin{grammarlens}[title={What you've built}]
One more word for telling, explaining and describing at B1.
\end{grammarlens}

\subsection*{Your turn}
\begin{itemize}
\raggedright
  \item \emph{Say it:} \emph{torris, torris}
  \item \emph{Say it:} \emph{torris, torris}, once more
  \item \emph{Say it:} a sentence of your own with \emph{torris, torris}
  \item \emph{Recall:} say the Latin for heat, then the Latin for glowing ash, then say \emph{torris, torris} again
\end{itemize}

\begin{tcolorbox}[breakable,colback=teal!4,colframe=teal!35!black,title={Before you move on}]
What does torris, torris mean? (\textquotedblleft{}A firebrand, a burning log\textquotedblright{}.) Say it once more.
\end{tcolorbox}
\section[ārdor, ārdōris]{ārdor, ārdōris — burning heat; passion}

\label{lesson:LA-C317-ardor}



\begin{quote}
Before the new one: say the Latin for glowing ash, then the Latin for a firebrand.
\end{quote}

\subsection*{You'll want to know: ārdor, ārdōris}
\textbf{ārdor, ārdōris} — \textquotedblleft{}burning heat; passion\textquotedblright{}.

\begin{quote}
Ārdor flammārum mūrōs nigrōs fēcit.
\end{quote}

\begin{quote}
\emph{The heat of the flames blackened the walls.}
\end{quote}

Masculine, 3rd declension, from \emph{ārdeō}, 'to burn' (taught). Also a burning eagerness. English \emph{ardour}.

\begin{cousinweb}
\textbf{ārdor} is built on \emph{ārdeō}, \textquotedblleft{}to burn\textquotedblright{} (EtymDB): burning, in fire and in feeling. English \emph{ardour} comes from it through French.
\end{cousinweb}

\begin{grammarlens}[title={What you've built}]
One more word for telling, explaining and describing at B1.
\end{grammarlens}

\subsection*{Your turn}
\begin{itemize}
\raggedright
  \item \emph{Say it:} \emph{ārdor, ārdōris}
  \item \emph{Say it:} \emph{ārdor, ārdōris}, once more
  \item \emph{Say it:} a sentence of your own with \emph{ārdor, ārdōris}
  \item \emph{Recall:} say the Latin for glowing ash, then the Latin for a firebrand, then say \emph{ārdor, ārdōris} again
\end{itemize}

\begin{tcolorbox}[breakable,colback=teal!4,colframe=teal!35!black,title={Before you move on}]
What does ārdor, ārdōris mean? (\textquotedblleft{}Burning heat; passion\textquotedblright{}.) Say it once more.
\end{tcolorbox}
